\begin{scanpage}{0132}{115}
For the purpose of refining the solutions which were found previously, it is necessary to assume that in actuality the currents flow only on the body's surface, and that a wave travelling from one edge to the other will undergo a perturbation at the latter. The process of forming the fringing field when this occurs may be investigated in the following way. The edge wave propagated from one of the edges is diffracted by the other edges; the waves arising with this in turn are diffracted by adjacent edges, etc. In this chapter, we will investigate the case when the dimensions of the surface faces are so large in comparison with the wavelength that it is sufficient to limit oneself to considering the diffraction of only the primary edge waves. This phenomenon we shall call secondary diffraction.

In this chapter, secondary diffraction by an infinitely long strip (\S\ 20--\S\ 23) and by a circular disk (\S\ 24) is studied. The solution of these problems may be obtained by means of the principle of duality from the solution of the diffraction problems for an infinite slit and a circular hole in a flat, ideally conducting screen. In the latter case, the physical treatment of diffraction of edge waves is significantly simpler; it is exactly for this reason, therefore, that almost all diffraction studies of edge waves are related to holes in a plane screen. However, we will not take such a path, but we shall investigate a strip and a disk directly. This approach has the advantage that it is easily generalized to the case of three-dimensional bodies.
\sectiontitle{20}{Secondary Diffraction by a Strip. Formulation of the Problem.}
Let an infinitely thin, ideally conducting strip of width $2a$ and unlimited length be orientated in space as shown in Figure 45. A plane electromagnetic wave incident normal to the strip's edges is directed at an angle $\alpha$ to the plane $xoz$ and has the following form:
\[\mathbf E=\mathbf E_0e^{ik(x\cos\alpha+y\sin\alpha)},\quad\mathbf H=\mathbf H_0e^{ik(x\cos\alpha+y\sin\alpha)}.\tag{20.01}\]
\end{scanpage}
